% sec-03-methods.tex - Methods (full)
\section{Methods}

\subsection{Repository Based LLM and Input Curation}

The method places a finite set of appropriately sized, high-quality inputs in a
single repository directory so that a 1M-context LLM can process them efficiently
\cite{adoptionguide}. The inputs are the paper template
(\texttt{inputs/llm-adoption}), the bibliography (\texttt{inputs/references.bib}),
and the four research sources (\texttt{research/document-types} and
\texttt{research/industry-workflow}), each pre-chunked with its own README so the
next processing step executes optimally. Two practical limits are respected: total
input is kept under 5 MB, and individual README files are kept under roughly 25K
tokens to avoid per-file errors. Outputs are generated in sections (one
\texttt{.tex} per paper section) so that text and code form right-sized chunks for
downstream reuse. Figure 4 shows the architecture.

\begin{mermaidfig}
\node[mmin] (i1) at (0,0) {inputs/llm-adoption\\(paper template)};
\node[mmin] (i2) at (0,-1.7) {inputs/references.bib\\(author works)};
\node[mmin] (i3) at (0,-3.4) {research/*\\pre-chunked + READMEs};
\node[mmproc] (ctx) at (4.5,-1.7) {1M-context LLM\\total input $<$ 5 MB\\README $<$ 25K tokens};
\node[mmproc] (gen) at (9.0,-1.7) {sub-prompts produce\\sectioned outputs};
\node[mmgoal] (o1) at (1.7,-5.6) {mermaid};
\node[mmin] (o2) at (5.0,-5.6) {draft-paper};
\node[mmaccent] (o3) at (8.3,-5.6) {full-paper};
\node[mmgoal] (o4) at (11.6,-5.6) {final-paper};
\draw[mmedge] (i1) to[out=0,in=120] (ctx.north west);
\draw[mmedge] (i2)--(ctx);
\draw[mmedge] (i3) to[out=0,in=240] (ctx.south west);
\draw[mmedge] (ctx)--(gen);
\draw[mmedge] (gen.south) to[out=-90,in=90] (o2.north);
\draw[mmedge] (gen.south) to[out=-90,in=90] (o3.north);
\draw[mmedge] (o1)--(o2); \draw[mmedge] (o2)--(o3); \draw[mmedge] (o3)--(o4);
\end{mermaidfig}
\figcaption{Figure 4. Repository based LLM document-generation architecture.
Curated inputs feed a 1M-context model under two size limits; sub-prompts produce
sectioned outputs that become the four stage directories, each reusable as a
right-sized chunk for future builds.}

\subsection{Single Prompt, Sub-Prompt Schedule, and Stages}

A single master prompt (\texttt{prompts/prompt-paper.md}) drives the whole build.
\emph{Process A} reads that prompt and writes the four stage sub-prompts in
\texttt{sub-prompts/}; \emph{Process B} then executes them in order, growing the
work from Mermaid figures to a draft, a full, and a final LaTeX paper, with outputs
accumulating as project-scale context for the next stage \cite{adoptionguide}. The
schedule is adapted from the trial-protocol build
(\texttt{trial-protocol/sub-prompts}). Figure 5 shows the generate-then-execute
structure.

\begin{mermaidfig}
\node[mmgoal] (mp) at (3.6,0) {Master prompt\\generate then execute};
\node[mmin] (sp1) at (0,-2.0) {prompt-1-mermaid};
\node[mmin] (sp2) at (2.4,-2.0) {prompt-2-draft};
\node[mmin] (sp3) at (4.8,-2.0) {prompt-3-full};
\node[mmin] (sp4) at (7.2,-2.0) {prompt-4-final};
\node[mmproc] (st1) at (0,-4.2) {Stage 1\\mermaid};
\node[mmproc] (st2) at (2.4,-4.2) {Stage 2\\draft};
\node[mmproc] (st3) at (4.8,-4.2) {Stage 3\\full};
\node[mmgoal] (st4) at (7.2,-4.2) {Stage 4\\final};
\node[mmaccent] (acc) at (10.4,-3.1) {context\\accumulates\\across stages};
\draw[mmedge] (mp) to[out=180,in=90] (sp1.north);
\draw[mmedge] (mp)--(sp2); \draw[mmedge] (mp)--(sp3);
\draw[mmedge] (mp) to[out=0,in=90] (sp4.north);
\draw[mmedged] (sp1)--(st1); \draw[mmedged] (sp2)--(st2);
\draw[mmedged] (sp3)--(st3); \draw[mmedged] (sp4)--(st4);
\draw[mmedge] (st1)--(st2); \draw[mmedge] (st2)--(st3); \draw[mmedge] (st3)--(st4);
\draw[mmedged] (st3.east) to[out=0,in=-90] (acc.south);
\draw[mmedged] (acc.north) to[out=90,in=0] (st4.north east);
\end{mermaidfig}
\figcaption{Figure 5. Process A generates the four sub-prompts; Process B executes
them sequentially as Stages 1 to 4. Outputs accumulate as context that informs the
next stage, so the final paper is produced under one prompt.}

\subsection{The Phase 1 Document Landscape}

Phase 1 oncology trials depend on large documents before, during, and after the
study \cite{srcWorkflowGem, srcDocTypesGPT}. Before the trial: the initial IND
dossier, the Investigator's Brochure, the clinical trial protocol, and the informed
consent form (often 20 or more pages, written to a 6th-to-8th grade reading level).
During the trial: serious adverse event and SUSAR safety narratives on 7-day and
15-day clocks, the annual Development Safety Update Report, protocol amendments, and
dose-escalation meeting minutes. After the trial: tables, listings, and figures, the
clinical study report under ICH E3, manuscripts and congress abstracts, and lay
summaries. Figure 6 maps this landscape.

\begin{mermaidfig}
\node[mmaccent] (hub) at (4.4,-2.4) {Repository LLM\\document generation};
\node[mmin] (b1) at (0,0) {IND dossier};
\node[mmin] (b2) at (0,-1.5) {Investigator's Brochure};
\node[mmin] (b3) at (0,-3.0) {Clinical protocol};
\node[mmin] (b4) at (0,-4.5) {Informed consent (ICF)};
\node[mmproc] (d1) at (4.4,1.1) {SAE/SUSAR 7/15-day};
\node[mmproc] (d2) at (8.8,1.1) {DSUR (annual)};
\node[mmproc] (d3) at (8.8,-0.4) {Protocol amendments};
\node[mmproc] (d4) at (8.8,-1.9) {Dose-escalation minutes};
\node[mmgoal] (a1) at (8.8,-3.6) {TLFs};
\node[mmgoal] (a2) at (8.8,-5.1) {CSR (ICH E3)};
\node[mmgoal] (a3) at (4.4,-5.5) {Manuscripts / lay summaries};
\draw[mmedge] (b1) to[out=0,in=160] (hub.north west);
\draw[mmedge] (b2)--(hub); \draw[mmedge] (b3)--(hub);
\draw[mmedge] (b4) to[out=0,in=200] (hub.south west);
\draw[mmedge] (hub) to[out=80,in=180] (d1.west);
\draw[mmedge] (hub) to[out=20,in=180] (d2.west);
\draw[mmedge] (hub)--(d3); \draw[mmedge] (hub) to[out=-20,in=180] (d4.west);
\draw[mmedge] (hub) to[out=-40,in=180] (a1.west);
\draw[mmedge] (hub) to[out=-60,in=180] (a2.west);
\draw[mmedge] (hub) to[out=-110,in=90] (a3.north);
\end{mermaidfig}
\figcaption{Figure 6. The large Phase 1 oncology document landscape, grouped into
before-trial inputs (left), during-trial process documents (upper right), and
after-trial outputs (lower right), all generated by the repository LLM.}

\subsection{Document Decision-Gate Taxonomy}

Whether faster authoring speeds a trial depends on the kind of gate a document
controls \cite{srcDocTypesGPT}. A hard gate means the trial legally or ethically
cannot proceed; a protocol-defined gate means the trial's own rules block the next
cohort, arm, or adaptation; and a decision gate means progress is legally possible
but the sponsor will not invest without an internal or regulatory decision. Faster
creation speeds the trial only when the document is on the critical path. Figure 7
shows the taxonomy and Table 1 states it in full.

\begin{mermaidfig}
\node[mmdec] (q) at (4.6,0) {Is the document\\on the critical path?};
\node[mmgoal] (g1) at (0,-2.6) {Hard gate};
\node[mmproc] (g2) at (4.6,-2.6) {Protocol-defined gate};
\node[mmaccent] (g3) at (9.2,-2.6) {Decision gate};
\node[mmin] (e1) at (0,-4.7) {IND, IRB approval,\\clinical-hold response,\\material amendments};
\node[mmin] (e2) at (4.6,-4.7) {cohort-review pkg,\\interim SAP / DSMB};
\node[mmin] (e3) at (9.2,-4.7) {EOP2 briefing,\\go/no-go at lock};
\node[mmctx] (n) at (4.6,-6.6) {Faster authoring helps only on the critical path};
\draw[mmedge] (q) to[out=180,in=90] (g1.north);
\draw[mmedge] (q)--(g2);
\draw[mmedge] (q) to[out=0,in=90] (g3.north);
\draw[mmedge] (g1)--(e1); \draw[mmedge] (g2)--(e2); \draw[mmedge] (g3)--(e3);
\draw[mmedged] (e2)--(n);
\end{mermaidfig}
\figcaption{Figure 7. The hard / protocol-defined / decision gate taxonomy. The
schedule value of faster authoring is realized only for documents on the critical
path.}

\vspace{0.2em}
\noindent\begin{tabularx}{\textwidth}{L{2.6cm} L{4.4cm} Y}
\toprule
\textbf{Gate type} & \textbf{Meaning} & \textbf{Example documents and effect of faster authoring} \\
\midrule
Hard gate & The trial legally or ethically cannot proceed & Initial IND, IRB approval, clinical-hold response, and material amendments; faster authoring starts the review clock and all sites sooner but cannot shorten a fixed review period. \\
Protocol-defined gate & The trial's own rules block the next cohort, arm, or adaptation & Cohort-review packages and the interim-analysis statistical plan or DSMB charter; faster authoring shortens the pause after the data mature but not the observation window itself. \\
Decision gate & Legally possible, but the sponsor will not invest without a decision & The End-of-Phase 2 briefing book and the go/no-go assessment after database lock; faster authoring lets the sponsor request the meeting and begin scheduling sooner. \\
\bottomrule
\end{tabularx}
\figcaption{Table 1. The three document decision-gate types, their meaning, and the
effect of faster authoring (adapted from the document-types research sources).}

\subsection{Before, During, and After Phase 1: Data and Document Workflow}

Data collection in modern oncology trials runs through a digital ecosystem
\cite{srcWorkflowGem}. Before the trial, clinical research coordinators extract data
from electronic health records into an electronic data capture (EDC) system and
record baseline tumor measurements by RECIST. During the trial, electronic case
report forms capture dose, pharmacokinetic, and pharmacodynamic data; electronic
clinical outcome and patient-reported outcome devices capture symptoms; and adverse
events are graded by CTCAE, with dose-limiting toxicities captured in the 3+3
window. After the trial, data undergo source document verification by clinical
research associates before the database is locked. Figure 8 shows the pipeline.

\begin{mermaidfig}
\node[mmin] (e1) at (0,0) {EHR extraction\\(CRCs)};
\node[mmin] (e2) at (0,-1.7) {EDC entry\\(Rave, Oracle)};
\node[mmin] (e3) at (0,-3.4) {RECIST baseline};
\node[mmproc] (d1) at (4.6,0) {eCRF: dose, PK/PD};
\node[mmproc] (d2) at (4.6,-1.7) {eCOA / ePRO symptoms};
\node[mmproc] (d3) at (4.6,-3.4) {AE logs (CTCAE),\\DLT in 3+3 window};
\node[mmproc] (a1) at (9.2,-0.85) {SDV by CRAs};
\node[mmgoal] (a2) at (9.2,-2.55) {Database LOCK};
\node[mmaccent] (a3) at (12.6,-1.7) {locked data\\feeds documents};
\draw[mmedge] (e1)--(e2); \draw[mmedge] (e2)--(e3);
\draw[mmedge] (e1) to[out=0,in=180] (d1.west);
\draw[mmedge] (e2) to[out=0,in=180] (d2.west);
\draw[mmedge] (e3) to[out=0,in=180] (d3.west);
\draw[mmedge] (d1) to[out=0,in=150] (a1.west);
\draw[mmedge] (d2) to[out=0,in=180] (a1.west);
\draw[mmedge] (d3) to[out=0,in=180] (a2.west);
\draw[mmedge] (a1)--(a2);
\draw[mmedge] (a2)--(a3);
\end{mermaidfig}
\figcaption{Figure 8. The before / during / after data-collection pipeline, from
EHR extraction through EDC, eCRF, eCOA/ePRO, and CTCAE-graded adverse events to
source document verification and database lock, after which the locked data feed
document generation.}

Medical writers act as the project managers of document creation, working in an
electronic document management system such as Veeva Vault with controlled,
version-tracked authoring and standardized templates \cite{srcWorkflowGPT,
srcWorkflowGem}. Figures 9, 10, and 11 show the before, during, and after authoring
flows respectively.

\begin{mermaidfig}
\node[mmin] (i1) at (0,0) {PI clinical rationale};
\node[mmin] (i2) at (0,-1.5) {Biostatistician rules};
\node[mmin] (i3) at (0,-3.0) {TransCelerate template};
\node[mmproc] (w) at (4.5,-1.5) {Medical writers\\in eDMS (Veeva Vault)};
\node[mmproc] (o1) at (9.0,0) {IND application};
\node[mmproc] (o2) at (9.0,-1.5) {Protocol, IB};
\node[mmproc] (o3) at (9.0,-3.0) {ICF (6th-8th grade)};
\node[mmgoal] (g) at (9.0,-4.6) {IRB / FDA submit\\30-day FDA clock};
\draw[mmedge] (i1) to[out=0,in=150] (w.north west);
\draw[mmedge] (i2)--(w);
\draw[mmedge] (i3) to[out=0,in=210] (w.south west);
\draw[mmedge] (w) to[out=0,in=180] (o1.west);
\draw[mmedge] (w)--(o2);
\draw[mmedge] (w) to[out=0,in=180] (o3.west);
\draw[mmedge] (o3)--(g);
\end{mermaidfig}
\figcaption{Figure 9. Pre-trial authoring. PI, biostatistician, and template inputs
are compiled by medical writers in an eDMS into the IND, protocol, Investigator's
Brochure, and informed consent form, which start the 30-day FDA clock on
submission.}

\begin{mermaidfig}
\node[mmin] (t1) at (0,0) {SAE / SUSAR event};
\node[mmin] (t2) at (0,-1.6) {Emerging PK/PD};
\node[mmin] (t3) at (0,-3.2) {Cohort completion};
\node[mmproc] (p1) at (4.6,0) {Safety narrative\\7 / 15-day};
\node[mmproc] (p2) at (4.6,-1.6) {DSUR; amendments};
\node[mmproc] (p3) at (4.6,-3.2) {Dose-escalation minutes};
\node[mmgoal] (pi) at (9.2,-1.6) {PI: causality and\\scientific justification};
\node[mmaccent] (s) at (9.2,-3.6) {LLM synchronizes\\protocol + consent +\\site documents};
\draw[mmedge] (t1)--(p1); \draw[mmedge] (t2)--(p2); \draw[mmedge] (t3)--(p3);
\draw[mmedge] (p1) to[out=0,in=150] (pi.north west);
\draw[mmedge] (p2)--(pi);
\draw[mmedge] (p3) to[out=0,in=210] (pi.south west);
\draw[mmedge] (pi)--(s);
\end{mermaidfig}
\figcaption{Figure 10. During-trial authoring is driven by safety data: narratives
on 7-day and 15-day clocks, the DSUR, amendments, and escalation minutes, with the
PI assessing causality and the LLM synchronizing the protocol, consent, and site
documents.}

\begin{mermaidfig}
\node[mmgoal] (lk) at (0,-1.5) {Database LOCK};
\node[mmproc] (tlf) at (4.0,-1.5) {Biostatistics:\\TLFs};
\node[mmproc] (csr) at (8.0,-1.5) {Writers: narrative\\-> CSR (ICH E3)};
\node[mmdec] (pi) at (11.6,-1.5) {PI review\\and sign-off};
\node[mmin] (m1) at (8.0,0.5) {Manuscripts / abstracts};
\node[mmin] (m2) at (8.0,-3.5) {Lay summaries};
\draw[mmedge] (lk)--(tlf); \draw[mmedge] (tlf)--(csr); \draw[mmedge] (csr)--(pi);
\draw[mmedge] (csr) to[out=90,in=180] (m1.west);
\draw[mmedge] (csr) to[out=-90,in=180] (m2.west);
\end{mermaidfig}
\figcaption{Figure 11. After-trial authoring: from database lock, biostatisticians
generate tables, listings, and figures; writers draft the narrative into the CSR
under ICH E3; the PI reviews and signs off; and manuscripts and lay summaries
branch off.}

\subsection{Figures, Tables, and Formatting Method}

All figures use one five-step palette so that emphasis reads by fill and stroke
weight alone and each Mermaid figure reproduces 1:1 as a TikZ \texttt{mermaidfig}
(Figure 12). Body text stays black, as in the paper template; only figures carry
color. Tables are set to the body-text width with ragged-right
(\texttt{>\{\textbackslash raggedright\textbackslash arraybackslash\}}) columns so
no cell shows large interword gaps, following the formatting strategies in
\texttt{trial-protocol/final-protocol/publication}. The paper section architecture
maps each section to one \texttt{.tex} file (Figure 14), and figure grounding makes
each diagram traceable to a Mermaid source (Figure 15).

\begin{mermaidfig}
\node[mmgoal] (c1) at (0,0) {goal \#8B2E3F};
\node[mmproc] (c2) at (0,-1.5) {proc \#2F5D7C};
\node[mmaccent] (c3) at (0,-3.0) {accent \#D08770};
\node[mmin] (c4) at (4.4,0) {input \#BFD7EA};
\node[mmctx] (c5) at (4.4,-1.5) {ctx \#F4F7F9};
\node[mmdec] (c6) at (4.4,-3.0) {warn \#D9D9D9};
\node[mmproc] (t) at (8.8,-1.5) {TikZ mm* node\\style};
\node[mmgoal] (g) at (12.4,-1.5) {Identical figure\\in LaTeX};
\draw[mmedge] (c1) to[out=0,in=120] (t.north west);
\draw[mmedge] (c2)--(t); \draw[mmedge] (c3) to[out=0,in=210] (t.south west);
\draw[mmedge] (c4) to[out=0,in=110] (t.north);
\draw[mmedge] (c5)--(t); \draw[mmedge] (c6) to[out=0,in=250] (t.south);
\draw[mmedge] (t)--(g);
\end{mermaidfig}
\figcaption{Figure 12. The five-color palette (plus the gray decision fill) maps 1:1
from the Mermaid \texttt{classDef} classes to the TikZ \texttt{mm*} node styles, so
every figure looks the same in LaTeX as on GitHub.}

\begin{mermaidfig}
\node[mmproc] (main) at (4.4,0) {main.tex\\\textbackslash input assembles all};
\node[mmgoal] (s1) at (0,-2.0) {Abstract + Keywords};
\node[mmproc] (s2) at (3.0,-2.0) {Introduction};
\node[mmctx] (s3) at (6.0,-2.0) {Table of Contents};
\node[mmproc] (s4) at (9.0,-2.0) {Methods};
\node[mmproc] (s5) at (0,-3.8) {Results};
\node[mmproc] (s6) at (3.0,-3.8) {Discussion};
\node[mmproc] (s7) at (6.0,-3.8) {Limitations};
\node[mmgoal] (s8) at (9.0,-3.8) {Conclusions};
\node[mmctx] (s9) at (4.4,-5.4) {References + Back Matter (CC)};
\draw[mmedge] (main)--(s1); \draw[mmedge] (main)--(s2); \draw[mmedge] (main)--(s3); \draw[mmedge] (main)--(s4);
\draw[mmedge] (main) to[out=-120,in=90] (s5.north);
\draw[mmedge] (main) to[out=-90,in=90] (s6.north);
\draw[mmedge] (main) to[out=-70,in=90] (s7.north);
\draw[mmedge] (main) to[out=-40,in=90] (s8.north);
\draw[mmedge] (main) to[out=-100,in=90] (s9.north);
\end{mermaidfig}
\figcaption{Figure 14. The paper section architecture. Each PAPER FORMAT section is
one \texttt{sections/*.tex} file assembled by \texttt{main.tex}; the Table of
Contents is generated after the Introduction.}

\begin{mermaidfig}
\node[mmin] (m) at (0,0) {mermaid/fig-NN.md\\(Python Mermaid)};
\node[mmproc] (d) at (4.4,0) {nodes, edges, palette,\\quantitative data};
\node[mmproc] (t) at (8.8,0) {TikZ mermaidfig\\(mm* styles)};
\node[mmgoal] (l) at (12.6,0) {Identical figure\\in draft/full/final};
\node[mmdec] (v) at (6.6,-2.2) {Verify twice:\\no overlaps, looseness,\\box spacing};
\draw[mmedge] (m)--(d); \draw[mmedge] (d)--(t); \draw[mmedge] (t)--(l);
\draw[mmedged] (v) to[out=120,in=-90] (d.south);
\draw[mmedged] (v) to[out=60,in=-90] (t.south);
\end{mermaidfig}
\figcaption{Figure 15. Figure grounding: each Mermaid source is reproduced as a TikZ
\texttt{mermaidfig} of the same complexity, then verified twice for overlaps, arrow
looseness, and box spacing.}

\subsection{Real-Time Auto-Commit and Auto-PR Monitorability}

To keep the human in the loop, the build commits each file to GitHub the moment it
is generated and opens a single continuously updated pull request, so the author can
review artifacts as they appear without intervention (Figure 13). Commit discipline
follows the rules: one commit for each of \texttt{main.tex}, the style file, the
bibliography, and the README, and one commit per section \texttt{.tex}; the
second-to-last commit of each stage fixes all errors, and the last performs the
remaining repository updates.

\begin{mermaidfig}
\node[mmproc] (g1) at (0,0) {LLM generates file};
\node[mmproc] (g2) at (3.6,0) {git commit};
\node[mmproc] (g3) at (7.0,0) {push to branch};
\node[mmaccent] (g4) at (10.6,0) {single PR\\(continuously updated)};
\node[mmgoal] (g5) at (10.6,-2.2) {author reviews\\on GitHub};
\node[mmctx] (g6) at (3.6,-2.2) {one commit per file\\(Rules 6, 7, 8)};
\draw[mmedge] (g1)--(g2); \draw[mmedge] (g2)--(g3); \draw[mmedge] (g3)--(g4);
\draw[mmedge] (g4)--(g5);
\draw[mmedged] (g5) to[out=180,in=0] (g6.east);
\draw[mmedged] (g6) to[out=90,in=-90] (g1.south);
\end{mermaidfig}
\figcaption{Figure 13. The real-time auto-commit and auto-PR loop. Each generated
file is committed and pushed at once into one continuously updated pull request that
the author monitors, with one commit per file per the commit rules.}
